\documentclass[preprint,12pt]{elsarticle}
% Shared set-up of the journal versions (Elsevier elsarticle class). Loaded after \documentclass{elsarticle}.
\usepackage{amsmath,amssymb,bm,mathtools}
\usepackage{fontspec}
\setmainfont{TeX Gyre Termes}\setsansfont{TeX Gyre Heros}\setmonofont{Latin Modern Mono}[Scale=0.9]
\usepackage{graphicx,float,caption,subcaption,booktabs,longtable,array,multirow,tabularx,rotating}
\usepackage{xcolor,setspace,lineno,microtype,xurl,siunitx,enumitem,pdflscape}
\definecolor{lk}{RGB}{0,70,140}
\usepackage[colorlinks=true,linkcolor=lk,citecolor=lk,urlcolor=lk,breaklinks=true,pdfborder={0 0 0}]{hyperref}
\usepackage[nameinlink,capitalise]{cleveref}
\crefname{figure}{Fig.}{Figs.}\Crefname{figure}{Fig.}{Figs.}
\crefname{equation}{Eq.}{Eqs.}\Crefname{equation}{Eq.}{Eqs.}
\sisetup{detect-all,separate-uncertainty=false}
\captionsetup{font=small,labelfont=bf}
\setlength{\emergencystretch}{3em}
\newcommand{\figdir}{../figures}

\newcommand{\doihl}[1]{doi: \href{https://doi.org/#1}{\nolinkurl{#1}}}
\newcommand{\dd}{\mathrm{d}}
\newcommand{\ET}{\mathit{ET}}
\newcommand{\Ent}{\mathit{ET}_{\mathrm{nt}}}
\newcommand{\Ks}{K_{\mathrm{s}}}
\newcommand{\Se}{S_{\mathrm{e}}}
\newcommand{\TAW}{\mathit{TAW}}
\newcommand{\repourl}{https://github.com/halilunaziru73-creator/Coupled-Water-Fluxes-Pore-Scale-Dynamics-Precision-Irrigation-Mulching-Feedback-Loops}
\newcommand{\repo}{\href{\repourl}{repository}}
\newcommand{\dataurl}{\repourl/tree/main/Data}
\newcommand{\codeurl}{\repourl/tree/main/Manuscripts}
\newcommand{\softurl}{\repourl/tree/main/Software}
\newcommand{\datalink}{\href{\dataurl}{Data}}
\newcommand{\codelink}{\href{\codeurl}{Code}}
\newcommand{\softlink}{\href{\softurl}{Software}}
\newcommand{\manuallink}{\href{\manualurl}{Manual}}
\newcommand{\downloadlink}{\href{\exeurl}{SWAFL.exe}}
\newcommand{\manualurl}{https://github.com/halilunaziru73-creator/Coupled-Water-Fluxes-Pore-Scale-Dynamics-Precision-Irrigation-Mulching-Feedback-Loops/raw/main/Software/v1.1.0/SWAFL_Desktop/docs/manual_v1.1.0/SWAFL_Manual_v1.1.0.pdf}
\newcommand{\exeurl}{https://github.com/halilunaziru73-creator/Coupled-Water-Fluxes-Pore-Scale-Dynamics-Precision-Irrigation-Mulching-Feedback-Loops/releases/latest/download/SWAFL.exe}
\newcommand{\csurl}{https://github.com/halilunaziru73-creator/Coupled-Water-Fluxes-Pore-Scale-Dynamics-Precision-Irrigation-Mulching-Feedback-Loops/tree/main/Data/Comprehensive_Summary}
\newcommand{\cslink}{\href{\csurl}{\nolinkurl{Data/Comprehensive_Summary}}}
\newcommand{\csavail}{\par\smallskip\textbf{Full simulation outputs.} Every value of every log of every run (the annual water and nitrogen balances, the harvests and the crop state, year by year), one compressed table per experiment (\texttt{E00\_Baselines} to \texttt{E16\_Edge\_Cases}), the one-row-per-run summary of all 10\,190 runs (\texttt{00\_summary\_all\_simulations}), the per-experiment summary (\texttt{00\_SUMMARY\_ALL\_EXPERIMENTS}) and the script that made them (\texttt{summarize\_results3.py}) are in the folder \cslink{} of the repository (\url{\csurl}); their 20-year means reproduce the run table within 2\,\%.}

% the page number at the top right of every page, the first page too
\makeatletter
\def\ps@plain{\def\@oddhead{\hfil{\small\thepage}}\let\@evenhead\@oddhead\let\@oddfoot\@empty\let\@evenfoot\@empty}
\g@addto@macro\ps@pprintTitle{\def\@oddhead{\hfil{\small\thepage}}\let\@evenhead\@oddhead}
\pagestyle{plain}
\makeatother

\journal{SoftwareX}
\hypersetup{pdftitle={SWAFL Desktop: equation-traceable irrigation scheduling and design with a graphical workbench for the Daisy model},pdfauthor={Naziru Halilu}}
\begin{document}

\begin{frontmatter}
\title{SWAFL Desktop: equation-traceable irrigation scheduling and design with a graphical workbench for the Daisy model}

\author[aua,upna]{Naziru Halilu\corref{cor1}}
\ead{halilunaziru73@gmail.com}
\cortext[cor1]{Corresponding author. ORCID: 0009-0006-2324-9599.}
\affiliation[aua]{organization={Agricultural University of Athens},city={Athens},country={Greece}}
\affiliation[upna]{organization={Public University of Navarre},city={Pamplona},country={Spain}}

\begin{abstract}
Irrigation software seldom shows which relation produced a number, and process models such as Daisy reach a decision only through hand-made steps in which the origin of a number is lost. SWAFL Desktop is an open, offline application for irrigation scheduling and design that joins three computing engines (a meta-model of a 10\,190-run Daisy ensemble, a daily soil-water model and closed-form surface, drip and sprinkler design modules) with a graphical workbench that plans, runs and explains Daisy simulations. Every output is traced to a numbered equation with a provenance label, every Daisy result shown comes from a Daisy run, and every computing part is tested against a reference outside the software.
\end{abstract}

\begin{keyword}
irrigation decision support \sep irrigation design \sep Daisy model \sep graphical user interface \sep equation traceability \sep software verification
\end{keyword}
\end{frontmatter}

\section*{Code metadata}
\begin{table}[H]\centering\small
\begin{tabularx}{\textwidth}{lp{5.2cm}X}
\toprule
Nr. & Code metadata description & Metadata\\\midrule
C1 & Current code version & v1.1.0\\
C2 & Permanent link to code/repository used for this code version & \url{\repourl}\\
C3 & Permanent link to Reproducible Capsule & Not available; all checks re-run from \texttt{Manuscripts/reproducibility}\\
C4 & Legal Code License & MIT; Daisy is distributed under its own licence\\
C5 & Code versioning system used & git\\
C6 & Software code languages, tools, and services used & JavaScript (Electron, Node.js), HTML/CSS; Python for the verification scripts; Daisy 7.1.14\\
C7 & Compilation requirements, operating environments \& dependencies & Node.js and Electron (build); Leaflet and geotiff.js and Daisy 7.1.14 for Windows (bundled); runs on Windows, macOS, Linux and in a web browser\\
C8 & If available, link to developer documentation/manual & \url{\manualurl}\\
C9 & Support email for questions & halilunaziru73@gmail.com\\\bottomrule
\end{tabularx}
\end{table}

\section*{Software metadata}
\begin{table}[H]\centering\small
\begin{tabularx}{\textwidth}{lp{5.2cm}X}
\toprule
Nr. & Software metadata description & Metadata\\\midrule
S1 & Current software version & 1.1.0\\
S2 & Permanent link to executables of this version & \url{\exeurl}\\
S3 & Permanent link to Reproducible Capsule & Not available\\
S4 & Legal Software License & MIT (source); end-user licence for the compiled application\\
S5 & Computing platforms/Operating Systems & Microsoft Windows (portable executable with Daisy); macOS and Linux (Electron); any modern web browser (everything except running Daisy)\\
S6 & Installation requirements \& dependencies & None for the Windows executable (no installation, administrator rights or network)\\
S7 & If available, link to user manual & \url{\manualurl}\\
S8 & Support email for questions & halilunaziru73@gmail.com\\\bottomrule
\end{tabularx}
\end{table}

\linenumbers
\section{Motivation and significance}
Between soil-water science and a decision about when, how much and by which method to irrigate lies a layer of software: schedulers, sensor dashboards, design calculators and crop-water models \cite{Datta2023,Garcia2020,Jones2004}. Three weaknesses recur. Tools report numbers without saying which equation produced them or whether it comes from the literature or from an in-house heuristic; they are tested on cases chosen by their authors, with results seldom re-runnable by others \cite{Hutton2016,Wilkinson2016}; and screening emulators are used as absolute predictors \cite{Bennett2013,Moriasi2007}.

Process models such as Daisy \cite{Hansen1990,Abrahamsen2000,Hansen2012,Holbak2026}, HYDRUS \cite{Simunek2008}, DSSAT \cite{Jones2003}, APSIM \cite{Keating2003} and AquaCrop \cite{Steduto2009} are the scientific backbone, but a simulation is specified in a text file, run by a program and read from files. That suits a model developer and stops a person with a question about a field: how much of the water a mulch saves is lost to drainage, which trigger suits a sandy soil, what a smaller nitrogen dose does to leaching. A set of simulations that answers such a question is rarely recorded with the answer, and when a model result reaches a decision tool by hand, its origin is easily lost. Meanwhile, classical irrigation design (volume-balance advance in surface irrigation \cite{Walker1987,Walker1989}, lateral hydraulics in drip \cite{KellerKarmeli1974,Keller1990}, nozzle and pattern analysis for sprinklers \cite{Christiansen1942,Playan2005}) is published as closed-form equations with worked examples, a free and under-used source of verification data.

SWAFL Desktop addresses these problems under two rules. The \emph{rule of traceability}: every number leads to a numbered equation carrying a provenance label, and every claim of fidelity is backed by a re-runnable test against a reference outside the software. The \emph{rule of origin}: a layer shows a Daisy result only from a Daisy run. It is intended for irrigation engineers, agronomists, extension services, teachers and students who need transparent estimates of schedules and designs, and for researchers who want to run designed Daisy experiments without writing set-up files.

\section{Software description}
\subsection{Software architecture}
SWAFL Desktop is an Electron application written in JavaScript (\cref{fig:arch}). The Windows program is one portable executable containing the runtime, the application, Daisy 7.1.14 with its library and the 1\,154-page manual \cite{HaliluManual}; it needs no installation, administrator rights or network. The same code runs in a web browser. The user works in one window with a tree of layers, the drawing of the layer, a dashboard with the same five tabs for every layer, and a panel of project properties (soil, tillage, crop, mulch, weather, irrigation method, trigger, nitrogen and a working date from 1980 to 2100) that every layer shares.

\begin{sidewaysfigure}\centering
\includegraphics[width=0.97\textheight]{\figdir/p3/fig01_architecture_A4.pdf}
\caption{Schematic overview of the main processes in SWAFL Desktop. Boxes indicate storages and solid lines flows; dashed black lines are flows out of the system, dotted lines are flows of information and dashed grey lines exchange of information between modules. Dashed purple boxes with a triangle are the drop-down controls and driving variables of the Properties panel, and the drivers of soil hydraulic properties (SHP). Dashed red lines and numbered circles mark the seven feedback loops. Pink indicates water, orange the crop, light brown nitrogen and rust soil structure and hydraulic state; purple hexagons are processes and information (Engines~1 and 2, trigger, design layers) and peach boxes are design layers and results. The brown line is the soil surface, with groundwater hatched at the base.}
\label{fig:arch}
\end{sidewaysfigure}

Four computing parts stand behind the layers. \emph{Engine~1} is an annual meta-model that composes, by ratios and differences, the 20-year means of 10\,190 Daisy simulations in 17 designed experiments \cite{HaliluData} into evapotranspiration, soil evaporation, transpiration, irrigation, yield, drainage and nitrogen leaching for the chosen settings; it fits no parameter. \emph{Engine~2} is a daily 40-layer soil-water model that gives dates, depths and schedules and is scaled so that its 20-year means equal those of Engine~1. The \emph{design layers} solve volume-balance advance with Kostiakov--Lewis intake for border, furrow and basin irrigation, the drip lateral by iterating emitter discharge and Hazen--Williams friction to convergence, and sprinkler discharge, range, uniformity and spacing. The \emph{Daisy workbench} (\cref{fig:wb}) writes Daisy set-ups, runs Daisy and reads its logs back; its bridge (a log parser, a seasonal summariser, a soil map, a question map and one live object from which the Daisy layers read) enforces the rule of origin.

\begin{figure}[!htbp]\centering
\includegraphics[width=\textwidth]{\figdir/p4/fig01_workbench.pdf}
\caption{Modules of the Daisy workbench and the hand-over. The chain runs from the planner and the decision base through the experiments, Daisy and the parsers to the hand-over; the modules below the dotted line are used before Daisy is started.}\label{fig:wb}
\end{figure}

\subsection{Software functionalities}
\begin{itemize}[nosep]
\item \textbf{Scheduling}: sensor depth and suction trigger against 20 years of daily soil water; a trigger law that expresses the triggers of ten soils as one root-zone depletion; next irrigation, trigger and wilting dates for the working date.
\item \textbf{Design}: layout, advance, infiltrated depth, efficiencies and a design table for surface irrigation; lateral hydraulics and emission uniformity for drip; nozzle, spacing, uniformity, pipe network and hydraulic grade line for sprinklers; a search for the most efficient inputs.
\item \textbf{Exploration}: the fate of water saved by mulch, evaporation regimes, tillage and consolidation, roots, seven soil--atmosphere feedback loops, 20-season reliability and water and nitrogen budgets, a crop library of 48 crops, and a field map with Sentinel-2 indices and management zones.
\item \textbf{Daisy experiments}: scenarios from the Daisy library, 14 prepared experiments of the published ensemble as compact designs of 6 to 72 runs, and the user's own designs, each checked against a catalogue of 54 factors of what Daisy can vary; five factors scale the van~Genuchten--Mualem parameters of every horizon. The number of runs set by the user is a ceiling, kept by a seeded balanced (Latin-hypercube) choice when the combinations are more, and every run is kept file by file next to Daisy.
\item \textbf{Reading runs}: a summary of the range of every result and the effect of every input on every result, journal-style figures of every quantity of every log, mass balances, and the root cause of every failed run located at the line of the set-up.
\item \textbf{Model-development tools}: a set-up editor, a Richards soil column, Daisy response curves and calibration by sequential search, CMA-ES \cite{HansenOstermeier2001} or Bayesian optimisation \cite{Mockus1978}.
\item \textbf{Traceability}: each tile, chart and arrow names its equation; the manual labels every equation as literature, data-driven or engineered and opens at the section of the layer in use.
\end{itemize}

\section{Illustrative examples}
\subsection{Design layers against textbook examples}
Examples~6.1 (furrow) and 6.4 (border) of Walker and Skogerboe \cite{Walker1987} were recomputed. Sixteen published furrow quantities were reproduced with a median error of 0.08\,\% and a maximum of 1.81\,\% (\cref{fig:ws}), and all border quantities within 0.5\,\%. Thirteen advance times of the FAO Irrigation Paper~45 border examples \cite{Walker1989} were reproduced within $-0.2$ to $+0.5$\,\%. The drip lateral loss agrees with the Christiansen reduction factor to 0.01\,\%, and the rebuilt drip solver agrees with an independent solver to $3.2\times10^{-7}$ up to a head loss of 60\,\% of the inlet head, warning the user beyond 20\,\%. The sprinkler layer follows the textbook's decision logic, and its nozzle discharge and wetted diameter are within 1.3\,\% and 1\,\%. These tests exposed three defects of the first release, all corrected: a fixed advance exponent that made slow surface advance too fast (up to $-29$\,\%), a drip solver stopped after four iterations (errors up to 81\,\%), and a sprinkler range coefficient that under-predicted the wetted diameter by 14\,\%.

\begin{figure}[!htbp]\centering
\includegraphics[width=\textwidth]{\figdir/p3/fig03_ws_verification.pdf}
\caption{Verification against Example 6.1. (a) Advance time $t_L$, (b) application efficiency $E_a$ (published versus SWAFL; grey line is 1:1), (c) relative differences of the 16 quantities (blue: advance and required times; green: efficiencies).}
\label{fig:ws}
\end{figure}

\subsection{Engines against held-out simulations}
Engine~1 was queried at the factor levels of 1\,000 held-out Daisy runs (sensor depth $\times$ trigger $\times$ rain on a coarse loam). It ranked them correctly (Spearman $\rho=0.96$ for irrigation, 0.92 for drainage, 0.97 for soil evaporation), but its raw values carried a baseline offset between experiment families; after aligning the baseline, the normalised RMSE was 6.2\,\% for evapotranspiration and 18.9\,\% for irrigation (\cref{fig:e1}, \cref{tab:e1}). Engine~1 is therefore a rank-faithful screening emulator. Engine~2 closed the water balance to a median 0.1\,\% of inflow and matched the annual irrigation of Engine~1 within 10\,\% in 90 of 100 scenarios; the exceptions are coarse soils at dry triggers. A sweep of $10^5$ Engine~1 scenarios takes about 2\,s.

\begin{figure}[!htbp]\centering
\includegraphics[width=\textwidth]{\figdir/p3/fig05_engine1_heldout.pdf}
\caption{Engine 1 versus Daisy for the 1\,000 held-out E12 runs. Orange: as shipped; blue: re-anchored. (a) ET, (b) irrigation, (c) drainage, (d) yield, (e) transpiration, (f) soil evaporation, (g) nitrogen leached. Grey lines are 1:1.}
\label{fig:e1}
\end{figure}
\begin{table}[!htbp]\centering\small
\caption{Engine 1 against 1\,000 held-out E12 runs (coarse loam), off-axis cells. Raw: as shipped; re-anchored: baseline aligned. nRMSE: RMSE as \% of the observed mean; $\rho$: Spearman rank correlation.}
\label{tab:e1}
\begin{tabular}{lrrrrrrr}
\toprule
 & \multicolumn{3}{c}{Raw} & \multicolumn{3}{c}{Re-anchored} & \\\cmidrule(lr){2-4}\cmidrule(lr){5-7}
Variable & $R^2$ & RMSE & bias & $R^2$ & nRMSE (\%) & bias & $\rho$\\\midrule
ET (mm) & $-2.06$ & 69.5 & $+53.1$ & $-0.11$ & 6.2 & $+5.5$ & 0.84\\
Irrigation (mm) & 0.18 & 97.4 & $+71.4$ & 0.79 & 18.9 & $+8.3$ & 0.96\\
Drainage (mm) & 0.72 & 21.7 & $+14.6$ & 0.82 & 48.7 & $+6.2$ & 0.92\\
Soil evaporation (mm) & 0.71 & 15.4 & $-7.9$ & 0.71 & 10.1 & $-7.9$ & 0.97\\
Transpiration (mm) & $-1.01$ & 29.3 & $+21.4$ & 0.10 & 8.7 & $+4.0$ & 0.64\\
Yield (Mg\,ha$^{-1}$) & $-1.82$ & 1.51 & $-0.79$ & $-2.33$ & 27.9 & $-1.15$ & 0.70\\
N leached (kg\,ha$^{-1}$) & $-0.84$ & 18.5 & $-14.7$ & 0.33 & 34.6 & $-0.3$ & 0.64\\\bottomrule
\end{tabular}
\end{table}

\subsection{Running and explaining Daisy experiments}
The published study took 17.5 days of Daisy time (median 145\,s per run). The 14 compact designs of the workbench contain 368 runs that stand for 7\,990 runs of the study, about 15 hours instead of 361. They recover the marginal response curves with a median error of 13.2\,\% of their range, so they are for orientation, and a curve that matters is re-run with more levels where it bends, which is one setting (\cref{fig:wbsteps}). Daisy 7.1.14 ran the set-ups written by the workbench for the library scenarios, deficit irrigation and the soil-hydraulic factors; scaling the hydraulic parameters of a loam towards a faster-draining soil raised first-season percolation from 177 to 222\,mm and removed runoff, the expected direction.

\begin{figure}[!htbp]\centering
\begin{subfigure}{0.49\textwidth}\centering\includegraphics[width=\textwidth]{\figdir/shots_jpg/wb_step1_input.jpg}\caption{Step 1: my data and weather (manual p.~760).}\end{subfigure}\hfill
\begin{subfigure}{0.49\textwidth}\centering\includegraphics[width=\textwidth]{\figdir/shots_jpg/wb_step2_library.jpg}\caption{Step 2: library of soils, crops and fertilisers (p.~583).}\end{subfigure}\\[3pt]
\begin{subfigure}{0.49\textwidth}\centering\includegraphics[width=\textwidth]{\figdir/shots_jpg/wb_step4_scenarios.jpg}\caption{Step 4: scenarios, a question and its runs (p.~762).}\end{subfigure}\hfill
\begin{subfigure}{0.49\textwidth}\centering\includegraphics[width=\textwidth]{\figdir/shots_jpg/wb_step6_logplots.jpg}\caption{Step 6: log plots of the runs (p.~763).}\end{subfigure}\\[3pt]
\begin{subfigure}{0.49\textwidth}\centering\includegraphics[width=\textwidth]{\figdir/shots_jpg/wb_step9_layers.jpg}\caption{Step 9: to the SWAFL layers, which layer answers which question (p.~764).}\end{subfigure}\hfill
\begin{subfigure}{0.49\textwidth}\centering\includegraphics[width=\textwidth]{\figdir/shots_jpg/wb_tool_setup_editor.jpg}\caption{Tool: set-up editor with bracket check (p.~783).}\end{subfigure}\\[3pt]
\begin{subfigure}{0.49\textwidth}\centering\includegraphics[width=\textwidth]{\figdir/shots_jpg/wb_tool_calibrate.jpg}\caption{Tool: calibration (daisypy-optim) (p.~779).}\end{subfigure}\hfill
\begin{subfigure}{0.49\textwidth}\centering\includegraphics[width=\textwidth]{\figdir/shots_jpg/wb_tool_learn_files.jpg}\caption{Tool: learn the Daisy files, the grammar of a .dai file (p.~792).}\end{subfigure}
\caption{The guided path of the Daisy workbench and three of its tools (real screenshots of release 1.1.0). Each step carries its step number, a one-line ``what to do here'', the manual page and Back and Next buttons.}\label{fig:wbsteps}
\end{figure}

For an experiment the workbench produces the summary of effects, figures of every log and the diagnosis of failed runs (\cref{fig:engine}). The bridge reads the logs exactly: six bundled logs read by the bridge and independently by \texttt{pandas} differed by at most $3\times10^{-15}$, and the chemical budget closed within 0.02\,\% of the applied amount. The soil column differed from an independent reference by 0.3\,\% and 0.9\,\% in drainage for sand and loam. Comparing the denitrification response curve with 792 outputs stored by the Daisy project found a median error of 40\,\%, caused by rescaling inputs that Daisy had already rescaled; after correction all 792 rows matched to $3\times10^{-15}$.

\begin{figure}[!htbp]\centering
\begin{subfigure}{\textwidth}\centering\includegraphics[width=0.92\textwidth]{\figdir/shots_jpg/wb_summary_effects.jpg}\caption{The summary of an experiment: the effect of every input on every result (blue rises, red falls; in brackets the share of the mean) and the summary in words.}\end{subfigure}

\vspace{4pt}
\begin{subfigure}{\textwidth}\centering\includegraphics[width=0.92\textwidth]{\figdir/shots_jpg/wb_every_log_lai.jpg}\caption{One of the figures of every log: leaf area index through the year, one panel per water limit, rainfed and irrigated, the mean of the seasons with 10 to 90\,\% as the band.}\end{subfigure}
\caption{Outputs of the experiment engine drawn from Daisy runs (screen captures of the workbench).}\label{fig:engine}
\end{figure}

\section{Impact}
SWAFL Desktop turns a published process-model ensemble, the Daisy model itself and the classical design theory of three irrigation methods into one offline program whose every number can be audited. A reviewer can select a displayed quantity, find its equation in the manual, check it against its source and re-run the test that supports it, without access to the developer. This changes the question asked of decision software from ``is the tool right?'', which a field trial must answer, to which outputs rest on published theory reproduced against its textbook, which on simulated ratios that preserve the ranking of options, and which on engineered relations that should be confirmed locally. The verification found four defects that were invisible on the screen, which supports shipping a test against an outside reference with every computing part.

Three routes from Daisy to a decision become available in one program (\cref{tab:routes}): the published ensemble through Engine~1 in microseconds per scenario, compact experiments in hours, and the user's own runs in minutes each. The software supports research that needs designed simulation experiments (mulch savings, trigger choice, deficit irrigation, nitrogen and chemical fate), teaching of irrigation engineering and soil--plant--atmosphere modelling, and extension work where internet is unreliable. The bridge design is not specific to Daisy and can be reused for other process models. No field validation, agronomic benefit or user study is claimed; Daisy is bundled for Windows only, and the operating envelopes reported here define where field measurements would be most informative.

\begin{table}[!htbp]\centering\small
\caption{The three routes from Daisy to a decision.}\label{tab:routes}
\begin{tabularx}{\textwidth}{>{\raggedright\arraybackslash}p{3.0cm}XXX}
\toprule & Ensemble and Engine 1 & Compact experiment (workbench) & Own live run\\\midrule
Soil and weather & ten soils, one weather record, the design of the study & as chosen in the workbench (library or own soil, bundled or own weather) & own soil table and weather\\
Time & microseconds per scenario & about 0.2 to 4 hours per experiment & minutes per run\\
Needs Daisy & no & yes (Windows) & yes (Windows)\\
Answers & ranking of options, response curves, regime maps within the design & the shape of a response for the chosen conditions & the budget of one set-up, for the layers of the suite\\
Main limit & fixed design; offset of absolute levels & coarse levels (workbench) & one run, one set-up\\\bottomrule
\end{tabularx}
\end{table}

\section{Conclusions}
SWAFL Desktop 1.1.0 is an open, offline application for irrigation scheduling and design in which every output is tied to a numbered equation, a provenance label and a re-runnable test, and in which a graphical workbench plans, runs and explains Daisy simulations under the rule that a Daisy result is shown only from a Daisy run. Its design layers reproduce textbook examples within 2\,\%, its annual meta-model ranks held-out scenarios correctly, its daily model conserves mass, its bridge reads Daisy logs exactly, and the defects found by verification are corrected. Field validation is the next step.

\section*{CRediT authorship contribution statement}
\textbf{Naziru Halilu:} Conceptualization, Methodology, Software, Validation, Formal analysis, Data curation, Visualization, Writing -- original draft, Writing -- review \& editing.

\section*{Declaration of competing interest}
The author declares that he has no known competing financial interests or personal relationships that could have appeared to influence the work reported in this paper.

\section*{Funding}
This research did not receive any specific grant from funding agencies in the public, commercial, or not-for-profit sectors.

\section*{Data availability}
The simulation ensemble is archived on Zenodo (\doihl{10.5281/zenodo.23043626}); every log of every run is in the repository folder \cslink{}. The verification scripts are in \url{\codeurl}.

\section*{Acknowledgements}
The Daisy model is developed and maintained at the University of Copenhagen.

\section*{Supplementary material}
Supplementary material associated with this article is available online. \textbf{Supplementary Material S1} gives the complete text of the study in the order of this article, with every section, figure, table, equation and appendix (numbered S1, S2, \ldots), including the list of symbols, the coupled soil--water--atmosphere system and its variables, the derivations and error propagation of the engines, all result views, the field case, the equation ledger of the 134 equations and all screenshots. \textbf{Supplementary Material S2} gives the derivations, the data and the mechanisms behind every table and figure, and the map from every number to its script and data file.

\bibliographystyle{elsarticle-num}
\bibliography{../tex/refs}
\end{document}
